\documentclass[11pt,a4paper]{article}
\usepackage[utf8]{inputenc}
\usepackage[T1]{fontenc}
\usepackage{geometry}
\geometry{margin=1in}
\usepackage{hyperref}
\usepackage{enumitem}
\usepackage{booktabs}
\usepackage{longtable}
\usepackage{xcolor}
\usepackage{titlesec}
\usepackage{amsmath,amssymb}

\titleformat{\section}{\Large\bfseries\color{blue!70!black}}{}{0em}{}
\titleformat{\subsection}{\large\bfseries\color{blue!50!black}}{}{0em}{}
\titleformat{\subsubsection}{\normalsize\bfseries\color{blue!30!black}}{}{0em}{}

\title{\textbf{Replication Plan}\\[0.5em]
\large Deformation Behavior of Annealed Cu$_{64}$Zr$_{36}$ Metallic Glass via Molecular Dynamics Simulations\\[0.3em]
\normalsize OSTI ID: 1609039}
\author{Auto-generated Replication Blueprint}
\date{\today}

\begin{document}
\maketitle
\tableofcontents
\newpage

%---------------------------------------------------------------------
\section{Paper Overview}
%---------------------------------------------------------------------
This paper uses classical molecular dynamics (MD) simulations to study the deformation behavior of annealed Cu$_{64}$Zr$_{36}$ metallic glass. The workflow involves: (1) generating an amorphous Cu--Zr alloy via melt-quenching, (2) annealing the glass at different temperatures below $T_g$, (3) performing uniaxial tensile/compressive deformation, and (4) analyzing stress--strain curves, shear band formation, Voronoi polyhedral analysis of local structure, and free volume evolution. The study uses LAMMPS with the Finnis--Sinclair (EAM) interatomic potential for Cu--Zr.

%---------------------------------------------------------------------
\section{Detailed Workflow}
%---------------------------------------------------------------------

\subsection{Phase 1: Sample Preparation}
\begin{enumerate}[label=\textbf{1.\arabic*}]
  \item \textbf{Build initial configuration.}
    \begin{itemize}
      \item Create a random solid solution of 64\% Cu and 36\% Zr atoms in an FCC supercell.
      \item System size: $\sim$10,000--50,000 atoms (as specified in the paper).
    \end{itemize}
  \item \textbf{Melt-quench protocol.}
    \begin{itemize}
      \item Heat to 2000\,K (well above melting point) and equilibrate in NPT.
      \item Quench to target temperature (e.g., 300\,K) at specified cooling rate (e.g., $10^{10}$--$10^{12}$\,K/s).
      \item Use Nos\'{e}--Hoover thermostat/barostat.
    \end{itemize}
  \item \textbf{Annealing.}
    \begin{itemize}
      \item Anneal at various sub-$T_g$ temperatures for specified durations (e.g., 1--10\,ns).
      \item Multiple annealing temperatures to study structural relaxation effects.
    \end{itemize}
\end{enumerate}

\subsection{Phase 2: Deformation Simulations}
\begin{enumerate}[label=\textbf{2.\arabic*}]
  \item \textbf{Uniaxial tension/compression.}
    \begin{itemize}
      \item Apply strain at constant engineering strain rate (e.g., $10^8$--$10^9$\,s$^{-1}$).
      \item Use periodic boundary conditions in lateral directions.
      \item Record per-atom stress and strain tensors.
    \end{itemize}
  \item \textbf{Record stress--strain curves.}
    \begin{itemize}
      \item Compute virial stress tensor; extract $\sigma_{zz}$ vs.\ $\varepsilon_{zz}$.
      \item Identify elastic modulus, yield stress, and ultimate strength.
    \end{itemize}
\end{enumerate}

\subsection{Phase 3: Structural Analysis}
\begin{enumerate}[label=\textbf{3.\arabic*}]
  \item \textbf{Voronoi tessellation analysis.}
    \begin{itemize}
      \item Compute Voronoi polyhedra indices $\langle n_3, n_4, n_5, n_6 \rangle$ for each atom.
      \item Track fractions of icosahedral ($\langle 0,0,12,0 \rangle$) and other key polyhedra.
    \end{itemize}
  \item \textbf{Free volume analysis.}
    \begin{itemize}
      \item Compute atomic volume from Voronoi cells.
      \item Track free volume evolution during annealing and deformation.
    \end{itemize}
  \item \textbf{Shear band visualization.}
    \begin{itemize}
      \item Compute local atomic shear strain ($\eta_i^{\text{Mises}}$) using OVITO.
      \item Visualize shear band formation during deformation.
    \end{itemize}
  \item \textbf{Radial distribution function} $g(r)$ to characterize short-range order.
\end{enumerate}

\subsection{Phase 4: Results Reproduction}
\begin{enumerate}[label=\textbf{4.\arabic*}]
  \item Reproduce stress--strain curves for different annealing conditions.
  \item Reproduce Voronoi polyhedra fraction evolution plots.
  \item Reproduce shear band snapshots.
  \item Reproduce $g(r)$ plots and structural order parameter trends.
\end{enumerate}

%---------------------------------------------------------------------
\section{Datasets}
%---------------------------------------------------------------------

\begin{longtable}{p{4.5cm} p{3.5cm} p{6cm}}
\toprule
\textbf{Dataset / Source} & \textbf{Type} & \textbf{Location / Notes} \\
\midrule
\endfirsthead
\toprule
\textbf{Dataset / Source} & \textbf{Type} & \textbf{Location / Notes} \\
\midrule
\endhead
Cu--Zr EAM potential & Interatomic potential & NIST Interatomic Potentials Repository \newline \url{https://www.ctcms.nist.gov/potentials/} \newline Specific potential: Mendelev et al.\ (2009) or as cited \\
\midrule
Initial atomic configuration & Generated & Created by LAMMPS \texttt{create\_atoms} with random alloy \\
\midrule
Simulation parameters & From paper & Cooling rate, annealing T/time, strain rate, system size \\
\bottomrule
\end{longtable}

%---------------------------------------------------------------------
\section{Models, Tools, and Software}
%---------------------------------------------------------------------

\subsection{Simulation Software}
\begin{longtable}{p{4.5cm} p{2.5cm} p{6.5cm}}
\toprule
\textbf{Tool / Library} & \textbf{Version} & \textbf{Purpose / URL} \\
\midrule
\endfirsthead
\toprule
\textbf{Tool / Library} & \textbf{Version} & \textbf{Purpose / URL} \\
\midrule
\endhead
LAMMPS & stable (2022+) & Molecular dynamics simulation engine \newline \url{https://www.lammps.org/} \\
\midrule
OVITO & $\geq 3.7$ & Visualization, Voronoi analysis, atomic strain \newline \url{https://www.ovito.org/} \\
\midrule
OVITO Python API & $\geq 3.7$ & Scripted analysis pipeline \newline \texttt{pip install ovito} \\
\midrule
Python (NumPy, SciPy) & $\geq 3.8$ & Post-processing, stress--strain curve analysis \\
\midrule
Matplotlib & $\geq 3.4$ & Plotting stress--strain, $g(r)$, Voronoi statistics \\
\midrule
Atomsk & latest & Optional: atomic structure manipulation \newline \url{https://atomsk.univ-lille.fr/} \\
\bottomrule
\end{longtable}

\subsection{Interatomic Potential}
\begin{itemize}
  \item \textbf{Cu--Zr EAM/Finnis--Sinclair potential}: Mendelev et al.\ (2009) or the specific potential cited in the paper. Available from the NIST Interatomic Potentials Repository.
\end{itemize}

\subsection{Hardware Requirements}
\begin{itemize}
  \item Multi-core workstation or small HPC cluster.
  \item GPU-accelerated LAMMPS (KOKKOS package) recommended for large systems.
  \item $\geq$ 16\,GB RAM; $\sim$100\,GB storage for trajectory data.
\end{itemize}

\end{document}
